\documentclass{article}
\usepackage{cancel}
\usepackage{gensymb}
\usepackage{tikz}
\usepackage{amsmath}
\usetikzlibrary{angles,quotes}
\usetikzlibrary{arrows.meta,calc}

\renewcommand{\thesubsubsection}{\alph{subsubsection})}
\def\hwNUM{9}

\title{Astro HW \hwNUM}
\author{Pierson Lipschultz}

\begin{document}
\maketitle
\setcounter{section}{\hwNUM}

\subsection{Extinction}

\subsubsection{Apparent Mag In front}
Given \dots
\[m-M=5\log(d)-5\]
\[M = -1.1 \]
\[d = 800\text{pc}\]
We find \dots
\[m = 5\log(d)-5 -1.1 \boxed{\approx 8.415}\]
\subsubsection{Apparent Mag Behind Nebula}
Given \dots
\[m-M=5\log(d)-5 + A_v\]
\[A_v = 1.1\]
\[d = 820\]
We find that \dots
\[m = 5\log(820)-5 -\cancel{1.1}+ \cancel{1.1}\boxed{ \approx 9.569}\]

\subsubsection{Error}
It would be further than it actually is. This is because the stars' brightness would be dimmer, which would incorrectly imply that the star is further away. 

Given \dots
\[m = M + 5\log(d)- 5 \]
\[10^{(m-M)/5} = d\]
We find that \dots
\[d = 10^{(9.569 +1.1 +5 )/5} \approx 1360.8\text{pc}\]

\[1360/820 \approx 1.65 \approx \boxed{65\% (\text{error})} \]


\subsection{Ionizing stars and HII regions}
\subsubsection{}
Given \dots
\[E = hv = hc/\lambda\]


\subsection{Conservation of angular momentum}
\subsubsection{Deriving 17.15}
Given \dots
\[\omega = 2/5 mrv\]
\[\omega_0 = 2/5 mr_0 v_0\]
\[\omega_f= 2/5 mr_f v_f\]

We find that \dots
\[2/5 mr_f v_f = 2/5 mr_0 v_0\]
\[\cancel{2/5 m}r_f v_f = \cancel{2/5 m}r_0 v_0\]
\[v_f = \frac{r_0}{r_f} v_0 \]

\subsubsection{Verify Scaling}
Given \dots
\[\frac{r_0}{r_f} v_0 = v_f \]
\[\frac{GM}{r_f^2} = \frac{v_f^2}{r_f}\]
% With the goal of \dots
% \[\boxed{r_f = \frac{v_0^2r_0^2}{GM}}\]
We find that\dots
\[\frac{GM}{r_f} = v_0^2 \frac{r_0^2}{r_f^2}\]
\[\frac{1}{\cancel{r_f}}r_f^{\cancel{2}} = v_0^2 \frac{r_0^2}{GM}\]
\[r_f = \frac{v_0^2 r^2 }{GM}\]

Given\dots
\[M = 2 \times 10^{30}\]
\[v = 100 \text{m/s}\]
\[r = 10^{10}\text{m}\]
\[G = 6.67 \times 10^{-11}\]
We find that\dots
\[r_f = \frac{100^2 \times 10^{10}}{6.67\times 10^{-11} \times {2 \times 10^{30}}}\]
\[r_f \approx \boxed{7.496 \times 10^{-15} \text{m}}\]


\subsubsection{Enhancement}

Given \dots
\[\omega_0 = 1/5 mr_0 v_0\]
\[\omega_f= 2/5 mr_f v_f\]

We find that \dots
\[2/5 mr_f v_f = 1/5 mr_0 v_0\]
\[2/\cancel{5} mr_f v_f = 1/\cancel{5} mr_0 v_0\]
\[\boxed{2\frac{r_0}{r_f} v_0 = v_f}\]

\subsection{Post-MS nucleosynthesis}
\subsubsection{Triple-alpha Energy Release}
Given \dots
\[E = mc^2\]
\[M_{He} = 4.0026032497 \text{u} \qquad M_c = 12\text{u}\]
\[u = 1.66053873 \times 10^{-27}\]
We find that \dots
\[\Delta M = 12\text{u} - 3 \times (M_{He} \times u) \approx -1.29685044 \times 10^{-29}\]
\[E_{rel} = (1.29685044 \times 10^{-29})(300000000)^2\]
\[E_{rel} \approx1.167165 \times 10^{-12}\]
\[E_{pNuc} = \frac{1.167165 \times 10^{-12}}{12} \boxed{\approx 9.7263783 \times 10^{-14} \text{J}}\]

\subsubsection{Star with He Core lifetime}
Given \dots
\[
M_{\text{s}} = 0.3 M_\odot = 0.3 \times 1.989 \times 10^{30} \approx 5.97 \times 10^{29}\ \mathrm{kg}
\]
\[
L = 25 L_\odot = 25 \times 3.828 \times 10^{26} \approx 9.57 \times 10^{27}\ \mathrm{W}
\]

\[E_{rxn} = 9.7263783 \times 10^{-14} \text{J}\]

\[
m_{\mathrm{He}} = M_{\mathrm{He}} \times u = 6.64648\times10^{-27}\ \mathrm{kg}
\]
We find that \dots

\[
N_{\mathrm{He}} = \frac{M_{\text{s}}}{m_{\mathrm{He}}} = 8.98\times10^{55}
\]

\[
N_{\mathrm{rxn}} = \frac{N_{\mathrm{He}}}{3} = 2.99\times10^{55}
\]

\[
E_{\mathrm{tot}} = N_{rxn} E_{rxn} = 3.49\times10^{43}\ \mathrm{J}
\]

\[
t_{He Burn} = \frac{E_{\mathrm{tot}}}{L}
   = \frac{3.49\times10^{43}}{9.57\times10^{27}}
   \boxed{\approx 3.64\times10^{15}\ \mathrm{s}
          \approx 1.16\times10^{8}\ \mathrm{yr}}
\]

\subsubsection{Energy released from C \(\rightarrow\) Si and Si \(\rightarrow\) Fe Fusion}

% Given \dots
% \[E = mc^2 \qquad c = 3 \times 10^8\]
% \[M_{He} = 4.0026032497 \text{u} \qquad M_C = 12\text{u} \qquad M_{Si} = 27.976926 \text{u} \qquad M_{Fe} = 55.934942\text{u}\]
% \[u = 1.66053873 \times 10^{-27}\]
% We find that \dots
% \[\Delta M_{C \rightarrow Si} = M_{Si} - (2(M_{C}) + M_{He})\]
% \[\Delta M_{C \rightarrow Si} = \left|27.976926 u - (2(12\text{u}) + 4.0026032497 \text{u})\right| \approx .0025677\text{u} \approx 4.2538\times 10^{-29}\text{kg}\]
% \[E_{C \rightarrow Si} = (\Delta M_{C \rightarrow Si}) c^2 = 4.2538\times 10^{-29} \times 3 \times (10^8)^2 \boxed{\approx 1.27914 \times 10^{-12}\text{J}}\]
% And \dots
% \[\Delta M_{Si \rightarrow Fe} = M_{Fe} - (2(M_{Si}))\]
% \[\Delta M_{Si \rightarrow Fe} = \left|55.934942 u - 2(27.976926\text{u})\right| \approx .01891\text{u} \approx 3.14\times 10^{-29}\text{kg}\]
% \[E_{Si \rightarrow Fe} = (\Delta M_{Si \rightarrow Fe}) c^2 = 3.14\times 10^{-29} \times 3 \times (10^8)^2 \boxed{\approx 9.42023 \times 10^{-13}\text{J}}\]


Given \dots
\[
E = mc^2,\qquad c = 3.00\times10^8\ \mathrm{m\,s^{-1}},\qquad
u = 1.66053873\times10^{-27}\ \mathrm{kg}
\]
\[
M_{\mathrm{He}} = 4.0026032497\ \mathrm{u},\quad
M_{\mathrm{C}} = 12.000000\ \mathrm{u},\quad
M_{\mathrm{Si}} = 27.976926\ \mathrm{u},\quad
M_{\mathrm{Fe}} = 55.934942\ \mathrm{u}.
\]

\paragraph{(i) \textbf{\(C \rightarrow Si\)}}

\[
^{12}\mathrm{C} + 4\,^{4}\mathrm{He} \rightarrow {}^{28}\mathrm{Si} + \gamma .
\]

\[
\Delta m_{C\to Si}
= \big(M_{\mathrm{C}} + 4 M_{\mathrm{He}}\big) - M_{\mathrm{Si}}
= 0.033486999\ \mathrm{u}.
\]

\[
E_{C\to Si} = \Delta m_{C\to Si}\,c^2
= \left(0.033486999\,u\right)\left(1.66053873\times10^{-27}\,\frac{\mathrm{kg}}{u}\right)
(3.00\times10^{8})^{2}
\approx 5.00\times10^{-12}\ \mathrm{J}.
\]

\[
E_n = \frac{E_{C\to Si}}{28} \approx 1.79\times10^{-13}\ \mathrm{J}.
\]

\paragraph{(ii) \(Si \rightarrow Fe\)}

\[
2\,^{28}\mathrm{Si} \rightarrow {}^{56}\mathrm{Fe} + \gamma .
\]
\[
\Delta m_{Si\to Fe} = \big(2 M_{\mathrm{Si}}\big) - M_{\mathrm{Fe}}
= 0.018910\ \mathrm{u}.
\]
\[
E_{Si\to Fe} = \Delta m_{Si\to Fe}\,c^2
= \left(0.018910\,u\right)\left(1.66053873\times10^{-27}\,\frac{\mathrm{kg}}{u}\right)
(3.00\times10^{8})^{2}
= 2.83\times10^{-12}\ \mathrm{J}.
\]


\[
E_n = \frac{E_{Si\to Fe}}{56} = 5.05\times10^{-14}\ \mathrm{J}
\]

\noindent
As expected, the energy released \emph{per nucleon} decreases as you go from C to Si to Fe.

\subsection{Extra Credit}
\subsubsection{Deriving \(v_{dm}\)}
Given\dots
\[Ke = 1/2 mv^2, \qquad GPE = -\frac{GMm}{r}\]
 Because the object starts at rest\dots
\[Ke - GPE_{f} = - GPE_0\]
We find that \dots
\[1/2mv_{dm}^2 - \frac{GMm}{r_0} = -\frac{GMm}{r}\]
\[1/2mv_{dm}^2 = -\frac{GMm}{r} + \frac{GMm}{r_0}\]
\[v_{dm}^2 = -2\left(\frac{GMm}{mr} - \frac{GMm}{mr_0}\right)\]
\[v_{dm}^2 = -2\left(\frac{GM\cancel{m}}{\cancel{m}r} - \frac{GM\cancel{m}}{\cancel{m}r_0}\right)\]

\[v_{dm} = -\sqrt{2GM\left(\frac{1}{r} - \frac{1}{r_0}\right)}\]

\subsubsection{Integrating \(v(r)\) as \(dr/dt\)}
Given \dots
\[M = (4\pi/3)R_0^3\rho_0\]
We find \dots
\[dr/dt = v_{dm} = -\sqrt{2GM\left(\frac{1}{r} - \frac{1}{r_0}\right)}\]
\[dt = \frac{dr}{-\sqrt{2GM\left(\frac{1}{r} - \frac{1}{r_0}\right)}}\]
\[t = \int_{R_0}^{r} \frac{1}{-\sqrt{2GM\left(\frac{1}{r} - \frac{1}{r_0}\right)}}dr\]
\[t = -\frac{1}{\sqrt{2GM}} \int_{R_0}^{r} \frac{1}{\sqrt{\left(\frac{1}{r} - \frac{1}{r_0}\right)}}dr\]

\[\qquad u = \frac{r}{r_0}, \qquad du = 1/r_0dr, \qquad r_0du=dr, \qquad R_0 \rightarrow 0, \qquad r \rightarrow 1 \]
\[t = -\frac{r_0}{\sqrt{2GM}} \int_{0}^{1} \frac{1}{\sqrt{\left(\frac{1}{ur_0} - \frac{1}{r_0}\right)}} du\]
\[t = -\frac{r_0 \sqrt{r_0}}{\sqrt{2GM}} \int_{0}^{1} \frac{1}{\sqrt{\left(\frac{1}{u} - \frac{1}{1}\right)}} du\]

\[t = -\frac{\sqrt{r_0^3}}{\sqrt{2GM}} \int_{0}^{1} \frac{1}{\sqrt{\left(\frac{1 - u}{u}\right)}} du\]

\[t = -\frac{\sqrt{r_0^3}}{\sqrt{2GM}} \int_{0}^{1}\sqrt{\left(\frac{u}{1 - u}\right)} du, \qquad \int_{0}^{1}\sqrt{\left(\frac{x}{1 - x}\right)} dx = \frac{\pi}{2}\]

\[t = -\left(\frac{\sqrt{r_0^3}}{\sqrt{2GM}}\right)\left(\frac{\pi}{2}\right)\]


\end{document}